\section{Introduction}
Shelf stocking couples ordering decisions with physical constraints. A placement
can satisfy its immediate position target while blocking access to a later
product. A reachable hand pose can still collide with another carton, and a
gripper closure command does not establish load-bearing contact. As free space
shrinks, a useful system must choose when to insert, slide, rearrange, switch
hands, reposition the base, or abandon an unsuccessful attempt.

Our target is continuous completion: cartons and bottles start on a nearby
table, the robot stocks the complete shelf, and earlier placements remain part
of the physical state. The eventual setting is a small house with an
omnidirectional bimanual platform. We start with the existing MABEL hands;
two- versus three-finger alternatives require measured comparisons. VR
teleoperation, recording, and replay will provide demonstrations and diagnostic
tools, but these components are not yet validated here.

The central research question is whether reusable task programs can choose
effective contact and recovery strategies across shelf instances when paired
with a controller whose feasibility estimates and execution outcomes are
explicit. Proposed contributions are a continuous stocking task family, an
auditable manipulation interface, and an evaluation of program synthesis and
controller choices. None is presented as an established performance advantage.
